\section{Pok\'emon Environment}\label{app:environment}

We run experiments on three Pok\'emon titles: Pok\'emon Red (Game Boy, 1996; we use the re-release compatible with the Game Boy Advance emulator), Pok\'emon Crystal (Game Boy Color, 2000), and Pok\'emon Emerald (Game Boy Advance, 2004). All three are single-player, turn-based role-playing games with long-horizon structure: overworld navigation, NPC dialogue, turn-based Pok\'emon battles, inventory management, and gated objectives (badges, plot milestones).

\paragraph{Interface.} The emulator exposes a screen buffer (rendered at the native resolution of each game: 160$\times$144 for Red/Crystal, 240$\times$160 for Emerald), which we upscale $2\times$ for the vision-language model, and a discrete button channel with eight inputs: \texttt{UP}, \texttt{DOWN}, \texttt{LEFT}, \texttt{RIGHT}, \texttt{A}, \texttt{B}, \texttt{START}, \texttt{SELECT}. Every observation step advances the emulator by a fixed number of frames (120) so that menu animations, battle text, and walking animations resolve between successive agent decisions.

\paragraph{Text map.} Because vision-language models have known difficulties with fine-grained spatial reasoning over pixel grids, we provide an ASCII text-map $m_t$ alongside the frame observation $o_t$ (\cref{sec:preliminaries}). The text map is derived from emulator memory and describes the visible tile grid around the player: walkable tiles (\texttt{.}), walls (\texttt{\#}), interactable tiles (\texttt{?}, e.g., signs and talkable objects), NPCs (\texttt{N}), ledges, and the player's position and facing. The map covers the current on-screen area plus a small margin of tiles just off-screen so the agent can plan one step ahead. The text map contains no walkthrough, no objective list, and no global map information beyond what is currently visible; it compensates for the VLM's limited spatial reasoning, not for domain knowledge.

\paragraph{Milestones and the button-press metric.} We use the milestone sequence from the PokeAgent Challenge~\citep{karten2026pokeagent}: a dense, canonical ordering of completion events that a run reaches roughly monotonically as it progresses through the game. Emerald has 31 canonical milestones through the completion of the 3rd gym, Red has 18 canonical milestones through the completion of the 3rd gym. The primary cost metric is cumulative \emph{button presses}, not tool calls: a single \texttt{press\_buttons} invocation emitting the list \texttt{[A,\,A,\,DOWN]} counts as three presses. This rewards compression in the action channel independent of how the harness structures its tool calls, and it makes $\mathcal{H}_{\min}$ (one button per step) directly comparable to harnesses that batch multi-step presses into one tool call.

\begin{figure}[t]
\centering
\includegraphics[width=\linewidth]{analysis/figures/emerald_full_speedrunning_route.pdf}
\caption{\textbf{Emerald Speedrunning Route.} Milestones from Littleroot Town (1) to aquiring the Dynamo Badge (31), with game frames from each waypoint. The geographic overview (right) maps key locations. This series of milestones require substantial exploration and backtracking; agents must navigate branching paths, and manage nonlinear dependencies between objectives. The current world record speedrun completes this segment in 1:00:57 minutes~\citep{keepingiticy2024emerald}.}
\label{fig:emerald_full_speedrunning_route}
\end{figure}

\begin{figure}[t]
\centering
\includegraphics[width=\linewidth]{analysis/figures/red_full_speedrunning_route.pdf}
\caption{\textbf{Pok\'emon Red Speedrunning Route.} Milestones from obtaining starter Pok\'emon (1) to aquiring the Thunder Badge (18), with game frames from each waypoint. The geographic overview (right) maps key locations. This series of milestones require substantial exploration and backtracking; agents must navigate branching paths, and manage nonlinear dependencies between objectives. The current world record speedrun completes this segment in 43:04 minutes; see \url{https://www.speedrun.com/pkmnredblue}.}
\footnotetext{website: https://www.speedrun.com/pkmnredblue}
\label{fig:red_full_speedrunning_route}
\end{figure}

\paragraph{Memory reader.} The PokeAgent emulator exposes a memory reader that reads structured game state (party, inventory, party HP, current map ID, dialogue text, battle status) from the emulator RAM. Tools such as the expert harness's A$^*$ pathfinder, battle type chart, and damage calculator use this reader. For $\mathcal{H}_{\min}$ and $\mathcal{H}_\mathrm{CH}$, the memory reader is exposed only through the text-map $m_t$ derivation and a handful of general primitives (e.g., \texttt{get\_party\_hp}); \emph{not} through pre-wired domain tools.
